% =========================================================
% Construction of the Real Numbers from Rational Approximation
% =========================================================

\section{Construction via Rational Approximation}

\section*{Motivation}

\begin{exposition}
The rational numbers $\mathbb{Q}$ form an ordered field but are
\emph{not complete}. There exist bounded sets of rationals with no
least upper bound in $\mathbb{Q}$, and there exist Cauchy sequences
of rationals that do not converge to a rational number.
\end{exposition}

\begin{exposition}
One way to construct the real numbers is therefore to enlarge
$\mathbb{Q}$ so that every coherent limiting process among
rationals corresponds to a number.
\end{exposition}

\begin{exposition}
The constructions of Cantor (Cauchy sequences) and Dedekind
(cuts) achieve this algebraically or order-theoretically.
A third viewpoint arises from \emph{rational approximation}.
\end{exposition}

\begin{exposition}
The idea is:
\end{exposition}

\begin{exposition}
\[
\mathbb{R}
=
\mathbb{Q}
\quad+\quad
\text{all infinite rational refinement processes}.
\]
\end{exposition}

\begin{exposition}
These refinement processes arise naturally from the
\emph{mediant construction} and the \emph{Stern--Brocot tree}.
\end{exposition}

\begin{exposition}
\bigskip
\end{exposition}

\subsection*{The Mediant Construction}

\begin{definitionbox}{Definition (Mediant)}

Let

\[
\frac{a}{b},\quad \frac{c}{d}
\]

be rational numbers with $b,d>0$.

The \textbf{mediant} is the rational number

\[
\frac{a+c}{b+d}.
\]

\end{definitionbox}

\begin{remark*}[Fully quantified form]

\[
\forall a,b,c,d\in\mathbb{Z},\ b,d>0:
\quad
\text{mediant}\!\left(\frac{a}{b},\frac{c}{d}\right)
=
\frac{a+c}{b+d}.
\]

\end{remark*}

\begin{proposition}[Mediant inequality]
\label{prop:stern-brocot-mediant-inequality}

If

\[
\frac{a}{b} < \frac{c}{d},
\]

then

\[
\frac{a}{b}
<
\frac{a+c}{b+d}
<
\frac{c}{d}.
\]

\hyperref[prf:stern-brocot-mediant-inequality]{Go to proof.}
\end{proposition}

\begin{remark*}[Fully quantified form]

\[
\forall a,b,c,d\in\mathbb{Z},\ b,d>0:
\left(
\frac{a}{b}<\frac{c}{d}
\Rightarrow
\frac{a}{b}
<
\frac{a+c}{b+d}
<
\frac{c}{d}
\right).
\]

\end{remark*}

\begin{remark*}[Intuition]

The mediant lies strictly between the two fractions.
It therefore provides a natural way to refine intervals
between rational numbers.

\end{remark*}

\begin{exposition}
\bigskip
\end{exposition}

\subsection*{The Stern--Brocot Tree}

\begin{theorem}[Stern--Brocot enumeration]
\label{thm:stern-brocot}

Every positive rational number appears exactly once in the
Stern--Brocot tree obtained by repeated mediant insertion.

\hyperref[prf:stern-brocot]{Go to proof.}
\end{theorem}

\begin{remark*}[Fully quantified form]

\[
\forall q\in\mathbb{Q}_{>0},
\quad
\exists ! \text{ finite path in the Stern--Brocot tree reaching } q.
\]

\end{remark*}

\begin{remark*}[Consequence]

Finite Stern--Brocot paths correspond exactly to rational numbers.

\end{remark*}

\begin{exposition}
\bigskip
\end{exposition}

\subsection*{Infinite Refinement and Irrational Limits}

\begin{exposition}
The Stern--Brocot process also produces infinite sequences of nested
intervals of rationals.
\end{exposition}

\begin{remark*}[Stern--Brocot limit theorem]
The limit behavior of infinite Stern--Brocot paths was established earlier as
\hyperref[thm:stern-brocot-limit]{Stern--Brocot Limit Theorem}. In the language
of nested rational intervals, the same idea says the following.

Let

\[
\left(\frac{a_n}{b_n}\right),\quad
\left(\frac{c_n}{d_n}\right)
\]

be sequences of rationals obtained by successive Stern--Brocot
refinement so that

\[
\frac{a_n}{b_n}
<
\frac{c_n}{d_n}
\]

and

\[
\left[\frac{a_{n+1}}{b_{n+1}},\frac{c_{n+1}}{d_{n+1}}\right]
\subseteq
\left[\frac{a_n}{b_n},\frac{c_n}{d_n}\right].
\]

Then there exists a unique real number

\[
x
=
\lim_{n\to\infty}\frac{a_n}{b_n}
=
\lim_{n\to\infty}\frac{c_n}{d_n}.
\]

\end{remark*}

\begin{remark*}[Fully quantified form]

\[
\forall
\text{ nested rational intervals } I_n,
\quad
\exists! x\in\mathbb{R}
\text{ such that }
x\in I_n
\text{ for all } n.
\]

\end{remark*}

\begin{remark*}[Interpretation]

\begin{itemize}
\item Finite Stern--Brocot paths correspond to rational numbers.
\item Infinite Stern--Brocot paths correspond to irrational numbers.
\end{itemize}

Thus every real number arises from an infinite refinement of
rational approximations.

\end{remark*}

\begin{exposition}
\bigskip
\end{exposition}

\subsection*{Real Numbers as Limits of Rational Processes}

\begin{definitionbox}{Definition (Real number via rational approximation)}
\begin{definition}[Real number via rational approximation]\label{def:real-limit-approx}

A \textbf{real number} may be defined as the limit of a convergent
sequence of rational numbers generated by a coherent approximation
process (for example, a Cauchy sequence or a nested sequence of
rational intervals).

\end{definition}
\end{definitionbox}

\begin{remark*}[Fully quantified form]

\[
\mathbb{R}
=
\{ (q_n)_{n\in\mathbb{N}} \subseteq \mathbb{Q}
\mid (q_n) \text{ is Cauchy} \} / \sim .
\]

\end{remark*}

\begin{remark*}[Conceptual summary]

The real numbers can be viewed as

\[
\boxed{
\mathbb{R}
=
\mathbb{Q}
\;\cup\;
\text{all limits of infinite rational approximations}.
}
\]

Different constructions implement this idea in different ways:

\begin{center}
\begin{tabular}{l l}
\toprule
Construction & Limiting process \\
\midrule
Cantor & Cauchy sequences of rationals \\
Dedekind & cuts in the ordered rationals \\
Stern--Brocot & infinite rational refinement paths \\
\bottomrule
\end{tabular}
\end{center}

All three constructions yield isomorphic complete ordered fields.

\end{remark*}
